%----------------------------------------------------------------------------------------
% RATIONALE VERGELIJKINGEN EN ONGELIJKHEDEN. Alle opgaven zijn eigen werk.
%
% Opbouw: Start, leerstof (theorie, werkwijze of uitleg, voorbeeld, korte oefening), en op een
% nieuwe pagina de Opdrachten (door elkaar, stijgende moeilijkheid), Koppel en Test jezelf.
% Labels van dit hoofdstuk beginnen met h3-.
%----------------------------------------------------------------------------------------

\mychapter{3}{Vergelijkingen en ongelijkheden}

\lkr{Leerplan B+S': LPD 18 en K1 (rationale functies), het deel over rationale vergelijkingen en ongelijkheden, ook in contexten. De twee klassieke fouten: een oplossing houden die de noemer nul maakt, en bij een ongelijkheid kruiselings vermenigvuldigen met een noemer waarvan het teken niet vastligt. Laat leerlingen beide fouten eens zelf maken en ontdekken (opdracht~\ref{oef:h3-o7}).}

Wanneer kost een stuk gemiddeld minder dan 4 euro? Hoe ver moet een voorwerp van een lens staan om een scherp beeld te geven? Zulke vragen geven vergelijkingen of ongelijkheden met de onbekende in een noemer. Je lost ze op met wat je over veeltermen en tekentabellen weet, als je op de noemers let.

Twee eigenschappen dragen het hoofdstuk. Bij een vergelijking mag je vermenigvuldigen met een noemer, zolang die niet nul is. Bij een ongelijkheid mag dat niet zomaar, maar heeft een breuk hetzelfde teken als het product van teller en noemer: dan helpt de tekentabel uit hoofdstuk 1.

\begin{oefeningen}[Start]
\oef{h3-start-1}

\oef{h3-verken}
\end{oefeningen}

%========================================================================================
\section{Rationale vergelijkingen}

\subsection{Vermenigvuldigen met de noemer}

Vermenigvuldig je beide leden met de noemer, dan verdwijnt de breuk en hou je een veeltermvergelijking over. Uit de 2de graad weet je dat je beide leden mag vermenigvuldigen met een getal dat niet nul is. Maar de noemer is geen vast getal: hij hangt af van $x$, en voor sommige $x$ is hij nul. Daarom schrijf je eerst op waar de noemers niet nul zijn.

\begin{theorie}
De \textbf{\textcolor{accent}{bestaansvoorwaarden}} van een vergelijking zeggen dat geen enkele noemer nul mag zijn.

\begin{stelling}{Vermenigvuldigen met een noemer}
Voor elke $x$ met $N(x) \neq 0$ geldt \quad $\formule{V(x) = W(x) \iff N(x) \cdot V(x) = N(x) \cdot W(x)}$
\end{stelling}

\begin{bewijs}
Neem een $x$ met $N(x) \neq 0$. Is $V(x) = W(x)$, dan vermenigvuldig je beide leden met hetzelfde getal $N(x)$. Is omgekeerd $N(x) \cdot V(x) = N(x) \cdot W(x)$, dan deel je beide leden door $N(x)$, en dat mag omdat $N(x) \neq 0$.
\end{bewijs}
\end{theorie}

Binnen de bestaansvoorwaarden heeft de veeltermvergelijking dus precies dezelfde oplossingen als de oorspronkelijke. Waar een noemer nul is, werkt het terugdelen niet: daar kan de veeltermvergelijking een oplossing hebben die de oorspronkelijke vergelijking niet heeft.

\voorbeeld{}{Los op: $\dfrac{1}{x} + \dfrac{1}{x - 1} = \dfrac32$.\label{vb:h3-vergelijking}}{%
\uitleg{Vermenigvuldig met het kleinste gemeenschappelijke veelvoud van de noemers, $2x(x - 1)$. Dan verdwijnen alle breuken tegelijk.}
\begin{stappenbreed}
x \neq 0,\ x \neq 1 & Bestaansvoorwaarden. \\
2(x - 1) + 2x = 3x(x - 1) & Vermenigvuldigen met $2x(x - 1)$. \\
3x^2 - 7x + 2 = 0 & Alles naar één lid. \\
x = \frac{7 \pm 5}{6} \iff x = 2 \text{ of } x = \tfrac13 & Discriminant $49 - 24 = 25$. \\
x \in \bigl\{\tfrac13, 2\bigr\} & Beide voldoen aan de bestaansvoorwaarden.
\end{stappenbreed}
Controle met $x = 2$: $\frac12 + \frac11 = \frac32$.}

\begin{oefeningen}
\needspace{6\baselineskip}
\oef{h3-kgv}
\end{oefeningen}

\subsection{Valse oplossingen}

Wat gebeurt er als de veeltermvergelijking een oplossing heeft die een noemer nul maakt? Die oplossing hoort er dan niet bij: in de oorspronkelijke vergelijking kun je ze niet eens invullen. Zo'n getal heet een valse oplossing.

\voorbeeld{}{Los op: $\dfrac{x^2}{x - 2} = \dfrac{4}{x - 2}$.\label{vb:h3-vals}}{%
\naast{%
\uitleg{Na vermenigvuldigen vind je ook een getal dat de noemer nul maakt: dat is geen oplossing.}
\begin{stappenvrij}{4.6cm}{4.2cm}
x \neq 2 & Bestaansvoorwaarde. \\[1ex]
x^2 = 4 & Vermenigvuldigen met $x - 2$. \\[1ex]
x = 2 \text{ of } x = -2 & \\[1ex]
x = -2 & $x = 2$ voldoet niet.
\end{stappenvrij}
De figuur toont $\frac{x^2}{x - 2} - \frac{4}{x - 2} = \frac{x^2 - 4}{x - 2}$: de rechte $y = x + 2$ met een gat in $x = 2$. Ze snijdt de $x$-as alleen in $-2$.}{%
\begin{tikzpicture}[x=.42cm,y=.42cm,font=\scriptsize]
  \rooster{-4}{1}{4}{-2}{1}{6}
  \kromme{-4}{4}{\x+2}
  \vol{(-2,0)} \open{(2,4)}
\end{tikzpicture}}}

\begin{oefeningen}
\needspace{6\baselineskip}
\oef{h3-vals}
\end{oefeningen}

\begin{werkwijze}
Schrijf de bestaansvoorwaarden op: geen enkele noemer mag nul zijn. Vermenigvuldig beide leden met het kgv van de noemers. Los de veeltermvergelijking op. Houd alleen de oplossingen die aan de bestaansvoorwaarden voldoen.
\end{werkwijze}

\begin{oefeningen}
\oef{h3-vergelijking}
\end{oefeningen}

%========================================================================================
\section{Rationale ongelijkheden}

\subsection{Het teken van een breuk}

Bij een ongelijkheid mag je niet zomaar met de noemer vermenigvuldigen: is die negatief, dan keert het teken om. In de Verken zag je wat er dan misgaat. Wie $\frac1x < 2$ vermenigvuldigt met $x$, vindt $x > \frac12$, en mist alle negatieve $x$: zo is $\frac{1}{-1} = -1 < 2$.

Breng daarom alles naar één lid en maak er één breuk van. Dan moet je alleen nog weten wanneer een breuk positief of negatief is, en daarvoor hoef je niet te vermenigvuldigen.

\begin{theorie}
\begin{stelling}{Teken van een breuk}
Voor elke $x$ met $B(x) \neq 0$ hebben \quad $\formule{\frac{A(x)}{B(x)}}$ \quad en \quad $\formule{A(x) \cdot B(x)}$ \quad hetzelfde teken.
\end{stelling}

\begin{bewijs}
Vermenigvuldig teller en noemer met $B(x) \neq 0$:
\bewijsregel{\frac{A(x)}{B(x)} = \frac{A(x) \cdot B(x)}{\bigl(B(x)\bigr)^2} .}
De noemer $\bigl(B(x)\bigr)^2$ is strikt positief, en delen door een positief getal verandert het teken niet.
\end{bewijs}
\end{theorie}

Een breuk is dus positief als teller en noemer hetzelfde teken hebben, en negatief als ze een ander teken hebben. In de tekentabel tellen de factoren van de noemer mee zoals die van de teller, met dezelfde regel als in hoofdstuk 1: bij een oneven multipliciteit wisselt het teken. Het enige verschil is dat de breuk in een nulwaarde van de noemer niet bestaat. Die duid je aan met een streep.

\voorbeeld{}{Los op: $\dfrac{x + 1}{x - 2} \geqslant 3$.\label{vb:h3-ongelijkheid}}{%
\naast{%
\uitleg{Trek $3$ af, zet op één noemer en stel de tekentabel op. Een nulwaarde van de teller hoort bij de oplossing (want $\geqslant$), een nulwaarde van de noemer nooit.}
\begin{stappenvrij}{4.5cm}{4.6cm}
\frac{x + 1 - 3(x - 2)}{x - 2} \geqslant 0 & Eén lid, één breuk. \\[2.5ex]
\frac{-2x + 7}{x - 2} \geqslant 0 & Nulwaarde $\frac72$, pool $2$. \\[1.5ex]
x \in \left]2, \tfrac72\right] & Uit de tekentabel.
\end{stappenvrij}
\figuren{\tekentabel{cccccc}{x & & 2 & & \frac72 & \\\hline -2x + 7 & + & + & + & 0 & - \\ x - 2 & - & 0 & + & + & +}{\text{breuk} & - & | & + & 0 & -}}}{%
\begin{tikzpicture}[x=.36cm,y=.36cm,font=\scriptsize]
  \rooster{-3}{1}{7}{-3}{1}{7}
  \draw[grijs,dashed] (2,-3) -- (2,7.3);
  \draw[accent,dashed] (-3,3) -- (7.3,3);
  \kromme{-3}{1.25}{(\x+1)/(\x-2)} \kromme{2.5}{7}{(\x+1)/(\x-2)}
  \vol{(3.5,3)}
\end{tikzpicture}}
De grafiek bevestigt het: ze ligt op of boven $y = 3$ tussen de asymptoot $x = 2$ en het snijpunt in $x = \frac72$.}

\begin{werkwijze}
Breng alle termen naar het linkerlid en zet ze op één noemer. Ontbind teller en noemer. Stel een tekentabel op met de nulwaarden van teller en noemer. Lees de oplossing af: de nulwaarden van de teller horen erbij bij $\leqslant$ of $\geqslant$, die van de noemer nooit.
\end{werkwijze}

\begin{oefeningen}
\needspace{6\baselineskip}
\oef{h3-ongelijkheid}
\end{oefeningen}

\subsection{Factoren met een even multipliciteit}

Na het ontbinden kan een factor meer dan eens voorkomen, in de teller of in de noemer. Zo'n factor met een even multipliciteit verandert het teken niet, maar het getal blijft in de tekentabel staan: als nulwaarde van de teller is de breuk er nul, als pool bestaat ze er niet.

\voorbeeld{}{Los op: $\dfrac{x^2 - x - 2}{x^2 - 6x + 9} \leqslant 0$.\label{vb:h3-dubbel}}{%
\uitleg{Ontbind teller en noemer. Begin de tekentabel rechts, en ga naar links: de dubbele pool verandert het teken niet.}
\begin{stappenbreed}
\frac{(x - 2)(x + 1)}{(x - 3)^2} \leqslant 0 & Som $1$, product $-2$; de noemer is een merkwaardig product. \\[2.5ex]
\text{nulwaarden } -1,\ 2;\quad \text{dubbele pool } 3 & Rechts van $3$ is de breuk positief. \\[1ex]
x \in [-1, 2] & De nulwaarden horen erbij, want $\leqslant$.
\end{stappenbreed}
\figuren{\tekentabel{cccccccc}{x & & -1 & & 2 & & 3 & }{\text{breuk} & + & 0 & - & 0 & + & | & +}}
Wie de noemer vergeet omdat hij \textit{toch altijd positief is}, schrijft $3$ niet in de tabel. Hier verandert dat niets aan de oplossing, maar bij $\geqslant 0$ zou je $3$ ten onrechte meenemen.}

\begin{oefeningen}
\needspace{6\baselineskip}
\oef{h3-dubbel}

\needspace{6\baselineskip}
\oef{h3-tekentabel}
\end{oefeningen}

\subsection{Twee breuken vergelijken}

Staan er in beide leden breuken, dan breng je ook alles naar één lid. Je kunt de ongelijkheid $f(x) \geqslant g(x)$ ook in een figuur lezen: ze geldt waar de grafiek van $f$ boven die van $g$ ligt, of ze snijdt.

\voorbeeld{}{Los op: $\dfrac{x}{x + 2} \geqslant \dfrac{1}{x}$.\label{vb:h3-tweebreuken}}{%
\naast{%
\uitleg{Trek $\frac1x$ af en zet op de noemer $x(x + 2)$. Vermenigvuldigen met $x$ of met $x + 2$ mag niet: hun teken ligt niet vast.}
\begin{stappenvrij}{4.6cm}{4.4cm}
\frac{x \cdot x - (x + 2)}{x(x + 2)} \geqslant 0 & Eén breuk. \\[2.5ex]
\frac{(x - 2)(x + 1)}{x(x + 2)} \geqslant 0 & Ontbind de teller. \\[2.5ex]
\text{nulwaarden } -1 \text{ en } 2 & Uit de teller. \\
\text{polen } -2 \text{ en } 0 & Alle enkelvoudig: het teken wisselt telkens.
\end{stappenvrij}
\figuren{\tekentabel{cccccccccc}{x & & -2 & & -1 & & 0 & & 2 & }{\text{breuk} & + & | & - & 0 & + & | & - & 0 & +}}
Oplossing: $x \in \left]-\infty, -2\right[ \cup \left[-1, 0\right[ \cup \left[2, +\infty\right[$.}{%
\begin{tikzpicture}[x=.36cm,y=.36cm,font=\scriptsize]
  \rooster{-6}{1}{5}{-4}{1}{5}
  \draw[grijs,thick] plot[domain=-6:-.25,samples=60,smooth] (\x,{1/\x}) plot[domain=.2:5,samples=60,smooth] (\x,{1/\x});
  \kromme{-6}{-2.5}{\x/(\x+2)} \kromme{-1.6}{5}{\x/(\x+2)}
  \vol{(-1,-1)} \vol{(2,.5)}
  \node[accent,above] at (4.3,1) {$f$};
  \node[grijs,right] at (.4,3.2) {$g$};
\end{tikzpicture}}
In de figuur ligt de grafiek van $f(x) = \frac{x}{x + 2}$ op of boven die van $g(x) = \frac1x$ in precies die drie delen; tussen $-2$ en $-1$ en tussen $0$ en $2$ ligt ze eronder.}

\begin{oefeningen}
\needspace{6\baselineskip}
\oef{h3-tweebreuken}

\needspace{10\baselineskip}
\oef{h3-grafisch}
\end{oefeningen}

%========================================================================================
\section{Toepassingen}

In een context heeft de onbekende vaak een betekenis die haar teken vastlegt: een aantal stuks, een afstand of een tijd is positief. Dan is ook de noemer positief, en mag je bij een ongelijkheid wel met de noemer vermenigvuldigen: vermenigvuldigen met een positief getal behoudt de ongelijkheid.

\voorbeeld{}{Bij een productie van $q$ stuks is de gemiddelde kost $G(q) = 2 + \dfrac{500}{q}$ euro. Vanaf hoeveel stuks is ze lager dan $4$ euro?\label{vb:h3-context}}{%
\naast{%
\uitleg{Hier is $q > 0$, dus je mag met $q$ vermenigvuldigen zonder dat het teken omkeert.}
\begin{stappenvrij}{3.6cm}{3.6cm}
2 + \frac{500}{q} < 4 & \\[1ex]
500 < 2q & Maal $q > 0$. \\
q > 250 & Vanaf $251$ stuks.
\end{stappenvrij}}{%
\begin{tikzpicture}[x=.009cm,y=.18cm,font=\scriptsize]
  \draw[black!14,very thin,xstep=100,ystep=2] (0,0) grid (600,12);
  \draw[->,grijs] (0,0) -- (640,0) node[right,text=black] {$q$};
  \draw[->,grijs] (0,0) -- (0,13) node[above,text=black] {$G$};
  \foreach \t in {200,400,600}{\node[below,inner sep=1.5pt] at (\t,0) {$\t$};}
  \foreach \t in {2,4,8,12}{\node[left,inner sep=1.5pt] at (0,\t) {$\t$};}
  \draw[grijs,dashed] (0,2) -- (600,2);
  \draw[accent,dashed] (0,4) -- (600,4);
  \kromme{50}{600}{2+500/\x}
  \vol{(250,4)}
\end{tikzpicture}}}

\begin{oefeningen}
\oef{h3-context}
\end{oefeningen}

\voorbeeld{}{Twee tuinslangen vullen samen een zwembadje in $12$ minuten. De dunne slang alleen doet er $10$ minuten langer over dan de dikke. Hoe lang doet elke slang er alleen over?\label{vb:h3-samen}}{%
\uitleg{Een oplossing kan wiskundig kloppen en toch geen betekenis hebben in de context. Kies de onbekende: de dikke slang heeft alleen $t$ minuten nodig, met $t > 0$. Per minuut vult ze dan $\frac1t$ van het badje. Samen vullen ze per minuut $\frac{1}{12}$.}
\begin{stappenbreed}
\frac1t + \frac{1}{t + 10} = \frac{1}{12} & Per minuut samen; $t \neq 0$, $t \neq -10$. \\[2.5ex]
12(t + 10) + 12t = t(t + 10) \iff t^2 - 14t - 120 = 0 & Maal $12t(t + 10)$. \\[1ex]
t = \frac{14 \pm 26}{2} \iff t = 20 \text{ of } t = -6 & Discriminant $676 = 26^2$. \\[2ex]
t = 20 & Een tijd is positief: $-6$ valt weg.
\end{stappenbreed}
De dikke slang doet er $20$ minuten over, de dunne $30$. Controle: $\frac{1}{20} + \frac{1}{30} = \frac{3 + 2}{60} = \frac{1}{12}$.}

\begin{werkwijze}
Kies de onbekende en zeg welke waarden ze kan aannemen. Stel de vergelijking of ongelijkheid op: tijd is afstand gedeeld door snelheid, werk per minuut is $1 / \text{tijd}$. Los op met de bestaansvoorwaarden, controleer in de context en antwoord in een zin.
\end{werkwijze}

\begin{oefeningen}
\oef{h3-werk}
\end{oefeningen}

%========================================================================================
\opdrachtenkop

\begin{oefeningen}
\needspace{14\baselineskip}
\oef{h3-o1}

\needspace{8\baselineskip}
\oef{h3-o9}

\needspace{14\baselineskip}
\oef{h3-o2}

\needspace{8\baselineskip}
\oef{h3-o13}

\needspace{12\baselineskip}
\oef{h3-o10}

\needspace{6\baselineskip}
\oef{h3-o14}

\needspace{8\baselineskip}
\oef{h3-o11}

\needspace{8\baselineskip}
\oef{h3-o12}

\needspace{8\baselineskip}
\oef{h3-o3}

\oef{h3-o16}

\needspace{9\baselineskip}
\oef{h3-o15}

\needspace{8\baselineskip}
\oef{h3-o19}

\oef{h3-o4}

\needspace{6\baselineskip}
\oef{h3-o20}

\needspace{8\baselineskip}
\oef{h3-o5}

\needspace{12\baselineskip}
\oef{h3-o18}

\oef{h3-o21}

\oef{h3-o6}

\needspace{12\baselineskip}
\oef{h3-o17}

\needspace{8\baselineskip}
\oef{h3-o22}

\needspace{8\baselineskip}
\oef{h3-o7}

\needspace{8\baselineskip}
\oef{h3-o24}

\needspace{9\baselineskip}
\oef{h3-o8}

\needspace{8\baselineskip}
\oef{h3-o23}

\needspace{8\baselineskip}
\oef{h3-o25}
\end{oefeningen}

%----------------------------------------------------------------------------------------
% KOPPEL: verbanden met andere leerstof
%----------------------------------------------------------------------------------------
\needspace{20\baselineskip}
\opdrachtsectie{Koppel}

\lkr{Kansrekenen en Rijen: kijk na wat al gezien is. Het pakket legt de volgorde niet vast.}

\begin{oefeningen}
\oef{h3-koppel-1}

\oef{h3-koppel-2}

\end{oefeningen}

%----------------------------------------------------------------------------------------
\opdrachtsectie{Test jezelf}

\begin{oefeningen}
\oef{h3-t1}

\oef{h3-t2}

\oef{h3-t3}

\oef{h3-t5}

\oef{h3-t4}

\oef{h3-t6}
\end{oefeningen}

\herhaalschema{%
  \fv{oef:h3-t1}{vb:h3-vergelijking}{oef:h3-kgv}%
  \fv{oef:h3-t2}{vb:h3-vals}{oef:h3-vals}%
  \fv{oef:h3-t3}{vb:h3-ongelijkheid}{oef:h3-ongelijkheid}%
  \fv{oef:h3-t5}{vb:h3-dubbel}{oef:h3-tekentabel}%
  \fv{oef:h3-t4}{vb:h3-context}{oef:h3-context}%
  \fv{oef:h3-t6}{vb:h3-samen}{oef:h3-werk}}

\ifzelfstudie
\section*{Maak je eigen samenvatting}
Sluit dit boek en schrijf in je schrift: waarom je met een noemer mag vermenigvuldigen binnen de bestaansvoorwaarden, de werkwijze voor rationale vergelijkingen, waarom je bij een ongelijkheid niet zomaar met de noemer vermenigvuldigt, het teken van een breuk met bewijs, de werkwijze met één breuk en een tekentabel, en de werkwijze voor een context. Vergelijk daarna met de kaders.
\fi
